\section{Data}
\label{sec:data}

\subsection{Corpus and task}

The corpus pools five public Roboflow Universe landmine-detection exports,
identified here by source id: \texttt{ds0}, \texttt{ds1}, \texttt{May\_lwir},
\texttt{Jan\_lwir}, and a small derived set folded into \texttt{ds0}. Each image
carries one class, \texttt{landmine}, and axis-aligned boxes. There is no
mine-type label; the task is presence and localisation of a visible
mine-like object. Table~\ref{tab:data} gives the frozen split.

\begin{table*}[t]
\centering
\small
\caption{Frozen dataset \texttt{landmine\_combined\_v2}. ``neg'' counts images
with an empty label file (deliberate negatives), not missing labels; there are
zero missing labels. Manifest digest
\texttt{748322d6\ldots c80de}.}
\label{tab:data}
\begin{tabular}{lrrrr}
\toprule
split & images & instances & neg & sources \\
\midrule
train & 3{,}907 & 11{,}550 & 20 & Jan 100, May 776, ds0 2{,}589, ds1 442 \\
valid & 181 & 419 & 7 & Jan 42, ds0 138, ds1 1 \\
test & 189 & 253 & 5 & Jan 10, ds0 151, ds1 28 \\
\bottomrule
\end{tabular}
\end{table*}

The class is heavily imbalanced in the detection sense and, more importantly,
heterogeneous in acquisition. Table~\ref{tab:shift} shows the source/split shift
that dominates the results later. \texttt{Jan\_lwir} images average $7.7$ boxes
per test image; \texttt{ds0} and \texttt{ds1} average about one. A model that
learns ``many objects'' from \texttt{Jan} and applies it to \texttt{ds0} will
generate false positives by construction, and vice versa. This is not a nuisance
to be averaged away: it is the principal source of variance in the test set, and
we return to it in Section~\ref{sec:persource}.

\begin{table*}[t]
\centering
\small
\caption{Source and split composition. The test set is dominated by
\texttt{ds0} images (151 of 189) but by \texttt{Jan} boxes (77 of 253).}
\label{tab:shift}
\begin{tabular}{llrrrr}
\toprule
split/source & & images & boxes & boxes/img & empty rate \\
\midrule
test & Jan\_lwir & 10 & 77 & 7.70 & 0.00\% \\
test & ds0 & 151 & 148 & 0.98 & 1.99\% \\
test & ds1 & 28 & 28 & 1.00 & 7.14\% \\
\midrule
train & Jan\_lwir & 100 & 740 & 7.40 & 0.00\% \\
train & May\_lwir & 776 & 7{,}706 & 9.93 & 0.00\% \\
train & ds0 & 2{,}589 & 2{,}643 & 1.02 & 0.77\% \\
train & ds1 & 442 & 461 & 1.04 & 0.00\% \\
\midrule
valid & Jan\_lwir & 42 & 287 & 6.83 & 0.00\% \\
valid & ds0 & 138 & 131 & 0.95 & 5.07\% \\
valid & ds1 & 1 & 1 & 1.00 & 0.00\% \\
\bottomrule
\end{tabular}
\end{table*}

\subsection{Two content defects, disclosed}

\paragraph{The \texttt{lwir} files are not thermal.} Sources named
\texttt{May\_lwir} and \texttt{Jan\_lwir} suggest long-wave infrared imagery. They
are not. Mean HSV saturation on a $24$-image sample is $0.38$, against $0.25$ for
the \texttt{ds0} RGB reference, and the images carry visible marker pins. They are
colour RGB photographs. We keep them, because removing them would drop the only
multi-object source and make the task artificially easy, and we disclose the
consequence: a model could exploit marker pins as a shortcut. A twelve-crop zoom
check around test mines found no consistent marker, paint, or flag cue, but we
cannot rule out a subtler one, and we do not make thermal claims anywhere in this
paper.

\paragraph{Two test labels are objectively wrong.} Manual review identified two
test boxes that are not on mines: \texttt{ds0\_0801} annotates a tree trunk and
\texttt{ds0\_0883} annotates a fence. Both were confirmed independently by a
Grad-CAM energy of exactly zero (Section~\ref{sec:gradcam}). We remove them from
ground truth and report results both with and without the correction
(Table~\ref{tab:labelfix}). No training labels were found to be wrong, so no
retraining was performed.

\begin{table}[t]
\centering
\small
\caption{Effect of removing two objectively wrong test boxes (no retraining).
Intervals in the source artifact are unchanged; the shift is a ground-truth
correction, not a model change.}
\label{tab:labelfix}
\begin{tabular}{lccc}
\toprule
arm & metric & before & after \\
\midrule
C (two-stage) & box recall & 0.9644 & 0.9721 \\
C (two-stage) & image recall & 0.9891 & 1.0000 \\
C (two-stage) & FP/image & 0.153 & 0.153 \\
A (YOLO deploy) & box recall & 0.9763 & 0.9841 \\
B (YOLO matched) & box recall & 0.9723 & 0.9801 \\
\bottomrule
\end{tabular}
\end{table}

\subsection{Provenance and redistribution}

The source archives were obtained as public exports, but their licence terms were
not recorded at download time and the original project URLs are no longer
available to this project. We therefore do not redistribute any image, label file,
or crop. What is released is the split manifest with a SHA-256 digest, the
removed-leak list with reasons, the audit CSVs and montages, the trained weights,
and the code. Anyone holding the original archives can reconstruct the exact
splits byte-for-byte from the manifest. Full detail is in
\texttt{DATA\_PROVENANCE.md}; the weights carry the upstream Ultralytics AGPL-3.0
terms (Section~\ref{sec:availability}). We regard the missing licence record as a
defect of this study and state it rather than inferring a permissive licence.
